\documentclass[11pt]{article}

\usepackage[a4paper,margin=1in]{geometry}
\usepackage[T1]{fontenc}
\usepackage[utf8]{inputenc}
\usepackage{lmodern}
\usepackage{microtype}
\usepackage{amsmath,amssymb,mathtools,bm}
\usepackage{tikz}
\usepackage[hidelinks,hypertexnames=false]{hyperref}
\usepackage{enumitem}

\setlength{\parindent}{1.5em}
\setlength{\parskip}{0pt}
\allowdisplaybreaks

\newcommand{\R}{\mathbb{R}}
\newcommand{\C}{\mathbb{C}}
\newcommand{\N}{\mathbb{N}}
\newcommand{\Z}{\mathbb{Z}}
\newcommand{\Sp}{\operatorname{Sp}}
\newcommand{\Trp}{\operatorname{Tr}^{+}}
\newcommand{\Ree}{\operatorname{Re}}
\newcommand{\Imm}{\operatorname{Im}}
\newcommand{\Argg}{\operatorname{Arg}}
\newcommand{\Circ}{\operatorname{Circ}}
\newcommand{\diag}{\operatorname{diag}}
\newcommand{\supp}{\operatorname{supp}}
\newcommand{\Lip}{\operatorname{Lip}}
\newcommand{\Sym}{\operatorname{Sym}}
\newcommand{\Ot}{\widetilde{\mathcal O}}
\newcommand{\OO}{\mathcal O}
\newcommand{\dd}{\,d}
\newcommand{\ii}{\mathrm{i}}
\newcommand{\e}{\mathrm{e}}
\newcommand{\id}{\mathrm{Id}}
\newcommand{\interior}[1]{\mathring{#1}}
\newcommand{\comp}[1]{\complement\!\left(#1\right)}

\title{Tunneling Effect for the Schrödinger Equation with a Magnetic Field}
\author{B. Helffer \and J. Sjöstrand}
\date{1987\\[1ex]\small English translation of the French manuscript}

\begin{document}
\maketitle

\begin{center}
\emph{Source: Annali della Scuola Normale Superiore di Pisa, Classe di Scienze, 4th series, vol. 14 (1987), no. 4, pp. 625--657.}
\end{center}

\medskip

We propose in this article to study the spectrum ``at the bottom of the well'' for a Schrödinger operator with magnetic field defined on $\R^n$,
\begin{equation}
P(h)=\sum_{j=1}^n(\ii h\partial_{x_j}-A_j)^2+V(x),
\tag{0.1}
\end{equation}
in the spirit of our first article [HE--SJ]$_1$, that is, in a semiclassical setting as $h\to0$.
For fixed $h>0$, the spectral theory for this type of operator has been studied in detail, in particular in a series of papers by J. Avron, I. Herbst and B. Simon [A--H--S] (see also the survey by Hunziker [HU] and the more recent work of A. Mohamed [MO]).

We assume here that $V(x)$ is a $C^\infty$ potential such that
\begin{equation}
\liminf_{|x|\to\infty}V(x)>0,
\qquad
V^{-1}(0)\neq\varnothing,
\tag{0.2}
\end{equation}
and we are interested in the spectrum in an interval $I(h)$ tending to $0$ (that is, $I(h)=(\alpha(h),\beta(h))$ with $\alpha(h)\to0$ and $\beta(h)\to0$ as $h\to0$).
Notice that $h$ need not run through the whole interval $(0,h_0]$; it may run only through a subset of $(0,h_0]$ having $0$ as an accumulation point.
We assume $A_j\in C^\infty$, and under hypothesis (0.2) it is classical that the spectrum in $I(h)$ is discrete.

For technical reasons, we shall only be able to treat the case where $A$ is small. We therefore introduce a real parameter $t$ and consider the operator
\begin{equation}
P_t(h)=\sum_{j=1}^n(\ii h\partial_{x_j}-tA_j)^2+V(x).
\tag{0.3}
\end{equation}
Our results will hold uniformly for $t\in[-t_0,t_0]$, with $t_0>0$ sufficiently small.

After studying the one-well problem in the spirit of [HE--SJ]$_1$, we shall study the splitting between the first two eigenvalues in the presence of a magnetic field and exhibit effects that are specific to the magnetic field.
Since the approach is close to the case $t=0$, we shall mainly emphasize the new points: decay of eigenfunctions, computation of the interaction, and we shall pass more quickly over other points.

The theory developed here certainly has applications to the semiclassical study of the Schrödinger equation with periodic potential and magnetic field. We hope to return to this question later.

\section{Decay of eigenfunctions and applications}

As in [HE--SJ]$_1$, the spectral analysis will depend on controlling eigenfunctions, associated with eigenvalues in $I(h)$, outside $V^{-1}(( -\infty,0])$. Let us first observe that Theorem 1.1 of [HE--SJ]$_1$ has the following immediate generalization.

\medskip
\noindent\textbf{Theorem 1.1.}
Let $\Omega\subset\R^n$ be a bounded domain with $C^2$ boundary. Let $V\in C^0(\overline\Omega,\R)$ and $\phi\in\Lip(\overline\Omega,\R)$. Then, for every $u\in C^2(\overline\Omega)$ with $u|_{\partial\Omega}=0$, one has
\begin{align}
&\int_\Omega \left|\nabla_{tA}^h\bigl(\e^{\phi/h}u\bigr)\right|^2\dd x
+\int_\Omega \bigl(V-|\nabla\phi|^2\bigr)\e^{2\phi/h}u\overline u\dd x
\notag\\
&\hspace{4cm}=\Ree\int_\Omega \e^{2\phi/h}(P_t(h)u)\overline u\dd x,
\tag{1.1}
\end{align}
where $P_t(h)$ is defined in (0.3) and
\[
\nabla_{tA}^h v=h\nabla v+\ii tAv.
\]

Theorem 1.1 has the same consequences as in the case without magnetic potential (Section 2 of [HE--SJ]$_1$), and we briefly recall the statements obtained in this way.
Assume therefore that
\begin{equation}
V^{-1}(( -\infty,0])=\bigcup_j U_j,
\tag{1.2}
\end{equation}
where the $U_j$ are pairwise disjoint compact sets of zero diameter for the Agmon distance $d_V$ associated with the metric $V_+(x)\,dx^2$.
Let $S>2S_0$, where
\begin{equation}
S_0=\inf_{k\ne j}d_V(U_j,U_k).
\tag{1.3}
\end{equation}
For every sufficiently small $\eta>0$, introduce
\begin{equation}
M_j^\eta=\interior{B}(U_j,S)\cap\comp{\bigcup_{k\ne j}B(U_k,\eta)},
\tag{1.4}
\end{equation}
where the balls are taken for the Agmon distance.
Let $I(h)$ be an interval as above, and suppose that there exists $a(h)>0$ such that
\begin{equation}
|\log a(h)|=o(1/h),\qquad h\to0.
\tag{1.5}
\end{equation}
Denote by $P_{t,M_j^\eta}(h)$ the Dirichlet realization in $M_j^\eta$. Denote by
\[
\lambda_1^t,\ldots,\lambda_{M^t}^t
\]
the eigenvalues of $P_t(h)$ contained in $I(h)$, by
\[
\mu_{j,1}^t,\ldots,\mu_{j,m_j^t}^t
\]
the eigenvalues of $P_{t,M_j^\eta}(h)$ contained in $I(h)$, and by
\[
u_i^t\in L^2(\R^n),\qquad \varphi_{j,k}^t\in L^2(M_j^\eta)
\]
corresponding orthonormal systems of eigenfunctions.

Recall that the number of eigenvalues satisfies bounds of the form
\begin{equation}
M^t+m_1^t+\cdots+m_N^t=\OO(h^{-N_0}),
\tag{1.6}
\end{equation}
where the constant in $\OO(\cdot)$ is uniform for $t\in[-t_0,t_0]$.
For the global problem this follows from estimate (1.4) of [C--S--S], itself deduced from Kato's inequality [KA]$_1$, which gives an $N_0$ such that the number of eigenvalues in $(-\infty,\lambda]$, for $\lambda<\liminf_{|x|\to\infty}V$, is bounded by $Ch^{-N_0}$ with $C$ independent of $t$.

As in [HE] or [HE--MA--RO], from [HE--RO] (when $A$ and $V$ have symbol-type properties at infinity) and from the eigenfunction decay results below, which only depend on (1.6), one obtains the following.

\medskip
\noindent\textbf{Theorem 1.2.}
Let $V\in C^\infty$ be semibounded and $A\in C^\infty$. Then, for every
$\lambda<\liminf_{|x|\to\infty}V$,
\[
N_h(P^t,\lambda)\stackrel{\mathrm{def}}{=}
\#\{\lambda_j^t(h):\lambda_j^t(h)<\lambda\}
=h^{-n}\left(\int_{\xi^2+V(x)<\lambda}\dd x\dd\xi+\OO(h)\right),
\]
where the $\OO(h)$ may be chosen uniformly for $t\in[-t_0,t_0]$, and $\lambda_j^t(h)$ denotes the sequence of eigenvalues of $P_t(h)$. Here $\lambda$ is assumed to be a noncritical value of $V$.

For Dirichlet problems one compares with the global problem and observes that $N_h(P^t_{M_j},\lambda)\le N_h(P^t,\lambda)$, which completes the proof of (1.6).
The result above is obtained either with a less good remainder and without uniform control in $t$ in [C--S--S], or under additional assumptions in [HE--RO]. An analogous theorem is announced by Ivrii [IV].

We now recall the eigenfunction decay results.

\medskip
\noindent\textbf{Lemma 1.3} (cf. [HE--SJ]$_1$).
Let $\phi_j(x)=d_V(x,U_j)$ for $x\in M_j$. Then, for every $\varepsilon>0$,
\begin{equation}
\left\|\nabla\!\left(\e^{\phi_j(x)/h}\varphi_{j,k}^t\right)\right\|_{L^2(M_j)}
+\left\|\e^{\phi_j(x)/h}\varphi_{j,k}^t\right\|_{L^2(M_j)}
=\OO(\e^{\varepsilon/h}),
\tag{1.7}
\end{equation}
uniformly for $t\in[-t_0,t_0]$.

\medskip
\noindent\textbf{Lemma 1.4} (cf. [HE--SJ]$_1$).
Let $\widetilde\phi(x)=\min_j d_V(x,U_j)$. Then, for every $\varepsilon>0$,
\begin{equation}
\left\|\nabla\!\left(\e^{(1-\varepsilon)\widetilde\phi(x)/h}u_k^t\right)\right\|_{L^2(\R^n)}
+\left\|\e^{(1-\varepsilon)\widetilde\phi(x)/h}u_k^t\right\|_{L^2(\R^n)}
=\OO(\e^{\varepsilon/h}),\qquad h\to0,
\tag{1.8}
\end{equation}
uniformly for $t\in[-t_0,t_0]$.

Assume now (as will be verified in our case) that
\begin{equation}
m_j^t(h)=m_j(h),
\tag{1.9}
\end{equation}
and that
\begin{equation}
\begin{gathered}
P_t(h),\;P_{t,M_j}(h)\quad\text{have no spectrum in the intervals}\\
[\alpha(h)-2a(h),\alpha(h))\qquad\text{and}\qquad
(\beta(h),\beta(h)+2a(h)].
\end{gathered}
\tag{1.10}
\end{equation}
Then one has the following.

\medskip
\noindent\textbf{Lemma 1.5} (cf. [HE--SJ]$_1$).
There exists a bijection $b(\eta,t,h)$ from
\[
\Sp P_t(h)\cap I(h)
\quad\text{onto}\quad
\bigcup_{j=1}^N\bigl(\Sp P_{t,M_j^\eta}(h)\cap I(h)\bigr)
\]
such that
\[
b(\eta,t,h)(\lambda)-\lambda
=\OO\!\left(\e^{-S_0/h+\varepsilon(\eta)/h}\right),
\]
where the $\OO$ is uniform with respect to $t$ and $\varepsilon(\eta)\to0$ as $\eta\to0$.

Up to this point everything looks strictly identical to the case without magnetic field. In the case of several wells, one expects tunneling effects to produce eigenvalue gaps of order $\OO(\e^{-S_0/h})$ when symmetries suggest, in the classical setting, multiple eigenvalues (for example a symmetric double well). The problem is that this is only an upper bound on the tunneling effect, which could be smaller, or even vanish. One of the first questions is therefore whether the eigenfunctions $\varphi_{jk}^t$---whose behavior along the minimal geodesics between the wells played an essential role when $t=0$---actually decay like $h^{-N_0}\e^{-\phi_j(x)/h}$, or whether they decay faster.

The model
\[
(\ii h\partial_{x_1}-t\tfrac{x_2}{2})^2
+(\ii h\partial_{x_2}+t\tfrac{x_1}{2})^2
+(x_1^2+x_2^2)
\]
has eigenfunctions of the form
\[
P_{t,k}\!\left(\frac{x}{\sqrt h}\right)
\exp\!\left[-\frac{\sqrt{1+t^2/4}}{h}\frac{x_1^2+x_2^2}{2}\right],
\]
where $y\mapsto P_{t,k}(y)$ is a polynomial, and shows that in some cases the decay is faster when $|t|\ne0$.

The techniques developed in this section do not seem capable of improving the decay control (except when $V$ is rotation invariant and, in dimension two, the magnetic field $B$ is constant). This will lead us to consider small $t$, in the hope of concluding by perturbation techniques in $t$. Thus, in order to obtain optimal results, we shall be forced either
\begin{itemize}[leftmargin=2em]
\item to assume analyticity of $V$ and $A$ near the $d_V$-minimal geodesics between the wells and to suppose $|t|<t_0$, with $t_0>0$ sufficiently small; or
\item without analyticity assumptions on $V$, to choose $t_0(h)\to0$ as $h\to0$, for instance of the size dictated by $-Ch\log(1/h)$ (see (2.40)).
\end{itemize}

\section{Constructions near a strict minimum of $V$}

Throughout this section we assume
\begin{equation}
V^{-1}(( -\infty,0])=\{x_0\},\qquad x_0=0,
\qquad V'(0)=0,\quad V''(0)\ \text{positive definite},
\tag{2.1}
\end{equation}
and
\begin{equation}
x_0\in\Omega,
\qquad
\liminf_{\substack{|x|\to\infty\\x\in\Omega}}V(x)>0,
\tag{2.2}
\end{equation}
where $\Omega$ is either an open set with $C^\infty$ boundary or $\R^n$.

The study of the spectrum at the bottom of the well is essentially analogous to that in Sections 3 and 4 of [HE--SJ]$_1$; we shall therefore indicate only the essential steps and the points that differ.
The operator under consideration now has symbol
\begin{equation}
p_t(x,\xi)=\sum_j(\xi_j-tA_j)^2+V(x).
\tag{2.3}
\end{equation}

\subsection*{Eikonal equation}

In [HE--SJ]$_1$, in order to understand more precisely what happens in the classically forbidden region, one introduces the symbol
\begin{equation}
q_t(x,\xi)=-p_t(x,\ii\xi)
=\sum_j(\xi_j+\ii tA_j)^2-V(x).
\tag{2.4}
\end{equation}
The new feature is that $q_t(x,\xi)$ is no longer real for $t\ne0$. This becomes a genuine problem as soon as $A$ is not the gradient of a real $C^\infty$ function.

In [HE--SJ]$_1$ there are essentially two steps. The first is formal (in powers of $x$); it localizes the spectrum modulo $\OO(h^\infty)$ and passes without serious difficulty to the case considered here. The second step is the construction of B.K.W. solutions in the $C^\infty$ or analytic category. Only the analytic construction will survive unchanged.

In what follows we work simultaneously either in the framework
\begin{equation}
\text{$V$ and $A$ are formal power series,}
\tag{2.5$_1$}
\end{equation}
or in the framework
\begin{equation}
\text{$V$ and $A$ are holomorphic in a neighborhood of $0$.}
\tag{2.5$_2$}
\end{equation}
By diagonalizing $V''(0)$ and, if necessary, changing gauge $A\mapsto A+\nabla\varphi$, we may assume that
\begin{equation}
V(x)=\sum_{j=1}^n\lambda_jx_j^2+\OO(|x|^3),
\tag{2.6}
\end{equation}
and
\begin{equation}
A_j(x)=\frac12\sum_{k=1}^n b_{jk}x_k+\OO(|x|^2),
\qquad b_{jk}\ \text{antisymmetric}.
\tag{2.7}
\end{equation}
Thus
\begin{equation}
q_t(x,\xi)=
\sum_j\left(\xi_j+\frac{\ii t}{2}\sum_kb_{jk}x_k\right)^2
-\sum_j\lambda_jx_j^2+\OO(|x,\xi|^3),
\tag{2.8}
\end{equation}
and the Hamilton vector field is
\begin{align}
H_{q_t}=2\Bigg[&\sum_j\left(\xi_j+\frac{\ii t}{2}\sum_kb_{jk}x_k\right)\partial_{x_j}
\notag\\
&+\sum_j\left\{\lambda_jx_j-
\sum_k\left(\xi_k+\frac{\ii t}{2}\sum_\ell b_{k\ell}x_\ell\right)
\frac{\ii t}{2}b_{kj}\right\}\partial_{\xi_j}\Bigg]
+\OO(|x,\xi|^2).
\tag{2.9}
\end{align}
The flow vanishes at $0$ and the fundamental matrix is
\begin{equation}
F_t=2
\begin{pmatrix}
-\dfrac{\ii tB}{2} & I\\[1ex]
\diag(\lambda_j)-\dfrac{t^2}{4}B^2 & -\dfrac{\ii tB}{2}
\end{pmatrix}.
\tag{2.10}
\end{equation}

For $t=0$ (cf. [HE--SJ]$_1$), the eigenvalues are $\pm2\sqrt{\lambda_j}$ and the eigenspaces corresponding to the positive (respectively negative) eigenvalues are
\begin{equation}
\Lambda_\pm^0:\qquad \xi_j=\pm\sqrt{\lambda_j}\,x_j,
\tag{2.11}
\end{equation}
which we write as $\xi=\pm\phi'(x)$ with
\[
\phi(x)=\frac12\sum_j\sqrt{\lambda_j}\,x_j^2.
\]
One can then construct two Lagrangian manifolds $\Lambda_\pm$ such that
\begin{equation}
T_{(0,0)}(\Lambda_\pm)=\Lambda_\pm^0,
\qquad H_q\ \text{is tangent to }\Lambda_\pm,
\tag{2.12}
\end{equation}
contained in $q_0^{-1}(0)$ and defined near $0$ by a holomorphic generating function $\psi_0$:
\begin{equation}
\Lambda_\pm:\qquad \xi=\pm\psi_0'(x),
\qquad x\in V\quad(V\ \text{a neighborhood of }0\text{ in }\C^n).
\tag{2.13}
\end{equation}
In this case $\psi_0$ is real for real $x$, and the construction also works in the $C^\infty$ category. For $t\ne0$ we lose this last property and, in the smooth category, will have to content ourselves with formal constructions.

\medskip
\noindent\textbf{Lemma 2.1.}
Let
\[
V(x)\sim\sum_{|\alpha|\ge2}V_\alpha x^\alpha
\]
(resp. let $V$ be holomorphic near $0$ with the same expansion) and assume (2.6). Let
\[
\vec A(x)\sim\sum_{|\alpha|\ge1}\vec a_\alpha x^\alpha
\]
(resp. let $A$ be holomorphic near $0$ with the same expansion) and assume (2.7). Then there exists $t_0>0$ such that, for $t\in[-t_0,t_0]$, there exists
\[
\psi_t(x)\sim\sum_{|\alpha|\ge2}\psi_{\alpha,t}x^\alpha
\]
(resp. a holomorphic $\psi_t$) satisfying
\begin{equation}
q_t(x,\psi_t'(x))=0,
\tag{2.14}
\end{equation}
\begin{equation}
\psi_t(x)|_{t=0}=\psi_0(x),
\tag{2.15}
\end{equation}
and
\begin{equation}
\Ree\psi_t''(0)\quad\text{is positive definite}.
\tag{2.16}
\end{equation}
If
\begin{equation}
\Lambda_+^t=\{\xi=\psi_t'(x)\},
\qquad
\Lambda_-^t=\left\{\xi=-\overline{\psi_t'(\overline x)}\right\},
\tag{2.17}
\end{equation}
then $\Lambda_+^t$ and $\Lambda_-^t$ are formal (resp. holomorphic) Lagrangian manifolds contained in $q_t^{-1}(0)$, and $H_{q_t}$ is tangent to both.

\medskip
\noindent\textbf{Remark 2.2.}
One has
\[
\overline{\psi_t'(\overline x)}=\psi_{-t}'(x),
\]
hence
\[
\Lambda_-^t=\{\xi=-\psi_{-t}'(x)\}.
\]

\noindent\textbf{Proof.}
In the formal case we refer to Section 1 of [HE--SJ]$_1$; in the holomorphic case the proof in the Appendix of [SJ]$_2$ adapts to the present situation (with holomorphic dependence on $t$). We spell out only the formal part of the computation that will be useful later.

\smallskip
\noindent\emph{a) The case homogeneous in $x$.}
This situation is extensively described in a family of papers devoted to hypoellipticity: J. Sjöstrand, L. Boutet de Monvel--F. Trèves, L. Boutet de Monvel--A. Grigis--B. Helffer, L. Hörmander, and others (see Hörmander's book [HO] and the references therein).

Writing $V=\diag(\lambda_j)$, observe first that $F_t/2$ has the same spectrum as the Hermitian matrices
\begin{equation}
\begin{pmatrix}
-\ii tB & \pm V^{1/2}\\
\pm V^{1/2} & 0
\end{pmatrix}.
\tag{2.18}
\end{equation}
Indeed it is enough to conjugate by
\[
U_\pm=
\begin{pmatrix}
I&0\\[0.5ex]
\mp\dfrac{t}{2}V^{-1/2}B&\pm\ii V^{-1/2}
\end{pmatrix}.
\]
It follows easily that the eigenvalues are real and have the form $\pm\mu_j(t)$, with
\begin{equation}
\begin{cases}
\mu_j(t)>0,\\
\mu_j(t)=\mu_j(-t),\\
\mu_j(0)=\sqrt{\lambda_j}.
\end{cases}
\tag{2.19}
\end{equation}
A small perturbation calculation based on (2.18) shows that the ``positive'' trace $\sum_j\mu_j(t)$, which determines the first eigenvalue of the operator approximating the bottom of the well to first order,
\[
\sum_j\left(\ii\partial_{x_j}-\frac t2\sum_kb_{jk}x_k\right)^2
+\sum_j\lambda_jx_j^2,
\]
satisfies
\begin{equation}
\Trp\!\left(\frac{F_t}{2}\right)
=\sum_j\sqrt{\lambda_j}
+\frac{t^2}{2}\sum_{j<k}
\frac{|b_{jk}|^2}{\sqrt{\lambda_j}+\sqrt{\lambda_k}}
+\OO(t^4).
\tag{2.20}
\end{equation}

We determine the quadratic part $\frac12Q^t$ of the phase $\psi_t(x)$ by looking for it as a power series in $t$:
\begin{equation}
Q^t=Q_0+\ii tQ_1+t^2Q_2+\cdots,
\tag{2.21}
\end{equation}
with $Q_1$ symmetric and
\begin{equation}
Q_0=\frac12\sum_j\sqrt{\lambda_j}\,x_j^2.
\tag{2.22}
\end{equation}
Equation (2.14) becomes
\begin{equation}
\left(Q_t-\frac{\ii t}{2}B\right)
\left(Q_t+\frac{\ii t}{2}B\right)=V.
\tag{2.23}
\end{equation}
Expanding in powers of $t$ gives
\begin{equation}
\begin{cases}
Q_0^2=V,\\
(Q_1-B/2)Q_0+Q_0(Q_1+B/2)=0,\\
-(Q_1-B/2)(Q_1+B/2)+Q_0Q_2+Q_2Q_0=0,\\
\qquad\vdots
\end{cases}
\tag{2.24}
\end{equation}
The existence of $Q_t$ for small $t$ follows from the implicit-function theorem applied to
\[
\Sym(n,\C)\times\R\ni(Q,t)\longmapsto f(Q,t)-V\in\Sym(n,\C),
\]
where
\[
f(Q,t)=\left(Q-\frac{\ii t}{2}B\right)
\left(Q+\frac{\ii t}{2}B\right).
\]
The tangent map in the $Q$ variable is
\[
A\longmapsto AQ_0+Q_0A,
\]
which is bijective on $\Sym(n,\C)$. Thus $Q^t$ depends holomorphically on $t$, and solving (2.24) gives
\begin{equation}
\begin{cases}
Q_0=V^{1/2}=\diag(\sqrt{\lambda_j}),\\[0.8ex]
(Q_1)_{ij}=\dfrac12b_{ij}
\dfrac{\sqrt{\lambda_i}-\sqrt{\lambda_j}}
{\sqrt{\lambda_i}+\sqrt{\lambda_j}},\\[1.2ex]
(Q_2)_{ij}=
\displaystyle\sum_\ell
\left(\frac{b_{i\ell}\sqrt{\lambda_i}}{\sqrt{\lambda_i}+\sqrt{\lambda_\ell}}\right)
\frac{1}{\sqrt{\lambda_i}+\sqrt{\lambda_j}}
\left(\frac{b_{\ell j}\sqrt{\lambda_j}}{\sqrt{\lambda_\ell}+\sqrt{\lambda_j}}\right),\\[0.8ex]
\qquad\vdots
\end{cases}
\tag{2.25}
\end{equation}
We note that $Q_2$ is positive, as expected from the eigenfunction decay estimates and from the fact that the eigenfunctions of
\[
\sum_j\left(\ii\partial_{x_j}-\frac t2\sum_kb_{jk}x_k\right)^2+\sum_j\lambda_jx_j^2
\]
are of the form $P(t,x)\e^{-Q^t(x)}$, with $P(t,x)$ a polynomial in $x$.

\smallskip
\noindent\emph{b) The nonhomogeneous case.}
We content ourselves with solving modulo $\OO(t^\infty)$. Lemma 2.1 says more, but we shall need this version in the $C^\infty$ case.
We seek $\psi_t$ in the form
\begin{equation}
\psi_t=\sum_k(\ii t)^k\psi_k,
\tag{2.26$_1$}
\end{equation}
where $\psi_0$ is the known solution from [HE--SJ]$_1$ of
\begin{equation}
q_0(x,\psi_0'(x))=0,
\qquad \psi_k\ \text{real},
\tag{2.26$_2$}
\end{equation}
and we normalize by
\begin{equation}
\psi_k(0)=0.
\tag{2.26$_3$}
\end{equation}
The equation is
\begin{equation}
(\nabla\psi_t+\ii tA)^2=V,
\tag{2.26$_4$}
\end{equation}
and identifying powers of $t$ gives
\begin{equation}
\sum_{\substack{k+j=\ell\\k,j\ge0}}\vec\Theta_k\cdot\vec\Theta_j
=
\begin{cases}
V,&\ell=0,\\
0,&\ell>0,
\end{cases}
\tag{2.27$_\ell$}
\end{equation}
where
\[
\vec\Theta_0=\nabla\psi_0,\qquad
\vec\Theta_1=\nabla\psi_1+A,\qquad
\vec\Theta_k=\nabla\psi_k\quad(k>1).
\]
For $\ell=1$ this reads
\[
2\nabla\psi_0\cdot\nabla\psi_1=-2\nabla\psi_0\cdot A,
\]
and for $\ell>1$,
\[
2\nabla\psi_0\cdot\nabla\psi_\ell
=-\sum_{\substack{k+j=\ell\\k,j>0}}\vec\Theta_k\cdot\vec\Theta_j.
\]
As was seen in [HE--SJ]$_1$, an equation of this type can be solved analytically (respectively in $C^\infty$) as soon as it can be solved formally, which causes no particular difficulty.

\medskip
\noindent\textbf{Remark 2.3.}
The preceding proof can be used when one imposes $|t|\le Ch^\delta$ with $\delta>0$, since it gives a solution of the eikonal equation modulo $\OO(h^\infty)$ even in the $C^\infty$ case.

The main conclusion of the discussion around Lemma 2.1 is the expansion
\begin{align}
\psi_t(x)
={}&\sum_j\sqrt{\lambda_j}\frac{x_j^2}{2}
+\frac{\ii t}{4}\sum_{i,j}
\frac{\sqrt{\lambda_i}-\sqrt{\lambda_j}}
{\sqrt{\lambda_i}+\sqrt{\lambda_j}}\,b_{ij}x_ix_j
\notag\\
&+\frac{t^2}{2}\sum_{\ell}\sum_{i,j}
\left(\frac{b_{i\ell}\sqrt{\lambda_i}\,x_i}{\sqrt{\lambda_i}+\sqrt{\lambda_\ell}}\right)
\frac{1}{\sqrt{\lambda_i}+\sqrt{\lambda_j}}
\left(\frac{b_{j\ell}\sqrt{\lambda_j}\,x_j}{\sqrt{\lambda_j}+\sqrt{\lambda_\ell}}\right)
\notag\\
&+\OO(t^3x^2)+\OO(x^3).
\tag{2.28}
\end{align}
The coefficient of $t^2$ is positive; hence
\begin{equation}
\Ree\psi_t(x)\ge \Ree\psi_0(x)-C(|t|+|x|)|x|^2.
\tag{2.29}
\end{equation}
This settles the problem for $t$ and $|x|$ sufficiently small. In fact, direct examination of the eikonal equation shows that
\begin{equation}
(\nabla\Ree\psi_t)^2
=V+(\nabla\Imm\psi_t+tA)^2.
\tag{2.30}
\end{equation}
Thus $\Ree\psi_t$ is, at least near $0$, the Agmon distance associated with the potential
\[
V+(\nabla\Imm\psi_t+tA)^2,
\]
and consequently
\begin{equation}
\Ree\psi_t\ge\Ree\psi_0.
\tag{2.31}
\end{equation}

\subsection*{B.K.W. construction}

We distinguish three cases:
\begin{itemize}[leftmargin=2em]
\item formal B.K.W.;
\item holomorphic B.K.W.;
\item $C^\infty$ B.K.W. with $|t|<Ch^\rho$, $\rho>0$.
\end{itemize}
To localize the eigenvalues modulo $\OO(h^\infty)$, the first construction is sufficient (under $C^\infty$ assumptions on $V$ and $A$). Holomorphy of $V$ and $A$ permits localization modulo $\OO(\e^{-\varepsilon_0/h})$. We shall also need to compare the B.K.W. constructions with true eigenfunctions, and the third construction can be useful. Notice in particular that if $|t|\le Ch$, then [HE--SJ]$_1$ applies as it stands even in the $C^\infty$ case.

Let us indicate how this connects with the theory of Section 3 of [HE--SJ]$_1$. First make a formal or holomorphic change of coordinates (Morse lemma) so that in the new coordinates
\[
\widetilde\psi_t(x)=\frac12x^2.
\]
In $\C^n\times\C^n$ this corresponds to a complex canonical transformation $\widetilde K_t$ sending the two Lagrangian manifolds associated with $\psi_t$ onto the Lagrangians $\xi=\pm\ii x$ in the new coordinates. We then use the technique of [HE--SJ]$_1$ by composing with the transformation $K_f$ introduced there (see the discussion around formula (3.11) in [HE--SJ]$_1$).

The canonical transformation $K_f\circ\widetilde K_t$ is quantized by the FBI transform
\[
(T_tu)(x,h)=h^{-n/2}C_n\int
\exp\!\left(\frac{\ii}{h}f(x,\theta_t(y))\right)b_t(x,y,h)u(y)\dd y,
\]
with
\[
f(x,z)=\frac{\ii}{2}(x-z)^2-\frac{\ii}{4}x^2,
\qquad
b_t(x,y,h)=1+\OO(t),
\]
where $b_t$ is an analytic (or formal) symbol to be determined and the integral is taken over a neighborhood $V$ of $0$ (or is interpreted formally by a formal stationary-phase theorem).
The canonical transformation sends $\Lambda_t^+$ to $\xi=0$, $\Lambda_t^-$ to $x=0$, and $\R^{2n}$ to an I-Lagrangian manifold determined by
\[
\xi=\frac{2}{\ii}\frac{\partial\phi_t}{\partial x},
\]
where $\phi_t$ is a small perturbation of $\phi_0$.
The FBI transform $T_t$ sends $S^{m,k}\e^{-\psi_t/h}$ to $S^{m,k}$ (by stationary-phase expansions).

We would now like to choose $b_t$ so that $T_t$ is formally unitary from $L^2(\R^n)$ onto germs of holomorphic functions at $0$, equipped with the scalar product of
$L^2(L(dx);\e^{-2\phi_t(x)/h})$.
For this purpose write $T_t=T_1^tA_t$, where $A_t$ is an analytic pseudodifferential operator and $T_1^t$ is obtained by taking $b_t(x,y,h)=1$.
Following the calculations of [HE--SJ]$_1$, p.~366, one obtains
\[
(T_1^tu_1\mid T_1^tu_2)_{\phi_t}
=h^{-3n/2}c\int
\overline{\widetilde u_2(\overline x,h)}\,
C_t(x,hD_x,h)\widetilde u_1(x,h)\dd x,
\]
where $C_t(x,hD_x,h)$ is a formally self-adjoint pseudodifferential operator close to the identity for $t$ sufficiently small.
Choosing $A_t$ self-adjoint with $A_t^2\equiv C_t^{-1}$, one obtains a unitary $T_t=T_1^tA_t$ (and, in the classical category, it may be rewritten in the integral form displayed above).

We have deliberately remained somewhat vague about which operations are merely formal. Everything is justified exactly as in [HE--SJ]$_1$ by Sjöstrand's theory [SJ]. With these remarks, Sections 3 and 4 of [HE--SJ]$_1$ remain valid in the formal setting under $C^\infty$ hypotheses, or in the analytic setting under analytic hypotheses.

We can now state the main consequences, whose proofs follow from those sections together with the preceding remarks. The first result comes from the quadratic approximation.

\medskip
\noindent\textbf{Proposition 2.4 (ground state).}
Assume that $V$ is $C^\infty$ and that hypotheses (2.1) and (2.2) hold. Then the Dirichlet realization in $\Omega$ of $P_t(h)$ has first eigenvalue
\begin{equation}
\lambda_t^1(h)=h\,\frac{\Trp F_t}{2}+\OO(h^{3/2}),
\tag{2.32}
\end{equation}
where $\Trp F_t/2$ is defined in (2.20), and the remainder is uniform for $t\in[-t_0,t_0]$.

\noindent\textbf{Proof.}
It is enough to carry out the calculation for the quadratic approximation. We have seen that
\[
\frac{\Trp F_t}{2}\ge\frac{\Trp F_0}{2}=\sum_j\sqrt{\lambda_j}.
\]
This agrees with the classical results of J. Avron, I. Herbst and B. Simon [A--H--S] and B. Simon [SI]$_1$, according to which the ground-state level rises when a magnetic field is introduced. The fact that comparison with the quadratic approximation makes an error of order $\OO(h^{3/2})$ is proved in [HE--SJ]$_1$ and in [SI]$_2$ (in fact one has $\OO(h^2)$).
We omit the analogous statements for excited levels (see Proposition 2.5).

\medskip
\noindent\textbf{Proposition 2.5.}
Under the hypotheses of Proposition 2.4, consider, for $t\in[-t_0,t_0]$, a simple eigenvalue $E_1^t$ depending continuously on $t$ of
\[
\sum_j\left(\ii\partial_{x_j}-\frac t2\sum_k b_{jk}x_k\right)^2
+\sum_j\lambda_jx_j^2.
\]
Then there exists a symbol $a_t(x,h)$ in $S^{m,k}$ (cf. [HE--SJ]$_1$, p.~364, with $\lambda=1/h$), depending holomorphically on $t$ for $|t|\le t_0$, and an eigenvalue $E^t(h)$ with expansion
\[
E^t(h)\sim hE_1^t+\sum_{j\ge2}h^jE_j^t,
\]
where $E_j^t$ depends holomorphically on $t$, such that
\begin{equation}
\bigl(P^t(h)-E^t(h)\bigr)
\bigl(a_t\e^{-\psi_t(x)/h}\bigr)\equiv0
\tag{2.33}
\end{equation}
in the sense of formal power series in $x$ and $h$.

\medskip
\noindent\textbf{Corollary 2.6.}
There exist $h_0>0$ and $t_0>0$ such that, for every $N$, there is $C_N>0$ with the following property: for every $h\in(0,h_0]$ and $t\in[-t_0,t_0]$, the interval
\[
[E_N^t(h)-C_Nh^N,\,E_N^t(h)+C_Nh^N],
\qquad
E_N^t(h)=\sum_{j\le N}h^jE_j^t,
\]
contains exactly one eigenvalue of $P_{t,\Omega}(h)$.

\noindent\textbf{Proof.}
The proof is identical to that of [HE--SJ]$_1$. The only difference is that $\psi_t(x)$ is merely a formal power series in $x$.

Thus the spectrum modulo $\OO(h^\infty)$ of the Dirichlet realization of $P_t(h)$ in $\Omega$ is determined, at least near the ground-state level.
It would be interesting to study the case where the corresponding eigenvalue of the quadratic model is multiple. This causes no new difficulty for fixed sufficiently small $t$, but we have not controlled uniformity in $t$ when the multiplicity changes near $t=0$.

Finally, since $a_t$ and $\psi_t$ are formal, it is not possible to describe the true eigenfunctions cleanly by their B.K.W. approximations as in the case $t=0$ in [HE--SJ]$_1$. We shall circumvent this difficulty either by assuming that $V$ and $A$ are analytic and $|t|<t_0$ with $t_0$ sufficiently small, or by imposing $|t|\le Ch^\rho$ with $\rho>0$ sufficiently large.

Let us first treat the analytic case. Assume that $V$ and $A$ are analytic in the regions where we wish to compare the true eigenfunction with its B.K.W. approximation. The function $\psi_t$ is holomorphic near $0$, but can be extended, as in [HE--SJ]$_1$, to neighborhoods of open minimal geodesics between $0$ and another well.
Then $a_t$ and $E^t$ are analytic symbols in $h$, depending holomorphically on $t\in B(0,t_0)$. If $\widetilde a_t$ and $\widetilde E_t$ denote realizations of these symbols (cf. [SJ]$_1$), we have the following.

\medskip
\noindent\textbf{Proposition 2.7 (analytic case)} (cf. Theorem 4.6 of [HE--SJ]$_1$).
With the preceding notation and the hypotheses of Proposition 2.5, there exist $\varepsilon_0>0$ and $t_0>0$ such that, for every $t\in[-t_0,t_0]$,
\begin{equation}
\bigl(P^t(h)-\widetilde E^t(h)\bigr)
\bigl(\widetilde a^t\e^{-\psi_t(x)/h}\bigr)
=\OO(\e^{-\varepsilon_0/h})\e^{-\Ree\psi_t(x)/h},
\tag{2.34}
\end{equation}
with a remainder uniform in $t$.

Taking account of Section 1, namely the $\e^{-d_V(x)/h}$ decay of eigenfunctions, we can improve these results. For $|t|\le t_0$ (after possibly decreasing $t_0$ and $\varepsilon_0$), (2.34) can be rewritten as
\begin{equation}
\bigl(P^t(h)-E^t(h)\bigr)
\bigl(a^t\e^{-\psi_t(x)/h}\bigr)
=\OO(\e^{-\varepsilon_0/h})\e^{-d_V(x)/h},
\tag{2.35}
\end{equation}
which gives the following consequence (cf. Theorem 5.8 of [HE--SJ]$_1$).

\medskip
\noindent\textbf{Corollary 2.8 (analytic case).}
Under the hypotheses of Proposition 2.7, an eigenfunction $u^t(h)$ of $P_{t,\Omega}(h)$ associated with the eigenvalue near $\widetilde E^t(h)$ in the sense of Corollary 2.6 can be normalized so that, in every open set $U$ with $\overline U\subset\Omega$ that is saturated by the $d_V$-minimal geodesics joining points of $U$ to $0$,
\begin{equation}
u^t(h)-h^{-n/4}\widetilde a^t\e^{-\psi_t/h}
=\OO\!\left(\e^{-d_V(x)/h-\varepsilon_0/h}\right),
\tag{2.36}
\end{equation}
with a remainder uniform in $t$.
If, in addition, one considers the ground-state level, then after possibly decreasing $t_0>0$, $\widetilde a^t$ is elliptic on $\overline U\times[-t_0,t_0]$, that is,
\[
\widetilde a^t\sim\sum_j a_j(t,x)h^j,
\]
and
\begin{equation}
a_0\ne0\qquad\text{on }\overline U\times[-t_0,t_0].
\tag{2.37}
\end{equation}

\medskip
\noindent\textbf{Remark 2.9.}
If $u^t(h)$ is an eigenfunction associated with $\lambda^t(h)$, then $\overline{u^t(h)}$ is associated with $\lambda^{-t}(h)$. Hence, for a simple eigenvalue,
\[
\lambda^t(h)=\lambda^{-t}(h).
\]

\medskip
\noindent\textbf{Corollary 2.10.}
For every $U$ of the preceding type there exists $t_0>0$ such that, for $t\in[-t_0,t_0]$,
\begin{equation}
u^t(h)(x)=\OO(h^{-N_0})\e^{-\Ree\psi_t/h}
\qquad\text{in }\overline U,
\tag{2.38}
\end{equation}
and analogous estimates hold for derivatives.

In view of (2.31), estimate (2.38) improves the result of Section 1 whenever the inequality in (2.31) is strict. We were not able to prove this estimate directly using the energy estimates of Section 1; analyticity was essential for this improvement. Estimates (2.36) and (2.38) will be very important in determining the entries of the interaction matrix between different wells.

In the $C^\infty$ case the situation is less favorable. We have nevertheless seen that the eikonal equation can be solved modulo $\OO(|t|^\infty)$. To construct a B.K.W. solution, the transport equations must also be solved modulo $\OO(|t|^\infty)$; this will be done below. If $|t|<Ch^\rho$ with $\rho>0$, this gives constructions modulo $\OO(h^\infty)$ after taking a realization of the phase $\psi_t$ by a Borel-type procedure.

A second condition needed for (2.38) is
\begin{equation}
\frac{\Ree\psi_t}{h}\le\frac{d_V}{h}-C\log h,
\tag{2.39}
\end{equation}
which, since the coefficient of $t$ in $\psi_t$ is imaginary, leads to
\begin{equation}
\frac{t^2}{h}\le -C\log h,
\qquad C>0.
\tag{2.40}
\end{equation}
Under this condition, the notions $\OO(h^\infty)\e^{-d_V(x)/h}$ and $\OO(h^\infty)\e^{-\Ree\psi_t/h}$ are equivalent, and the theory of [HE--SJ]$_1$ applies without serious difficulty.

Let us briefly study the transport equations. We work at the ground-state level and follow the approach described in [HE]. Keep $t$ as an independent parameter and expand
\[
a^t(x,h)=a_0^t(x)+ha_1^t+\cdots.
\]
We analyze only the first three equations, viewing the $a_i^t(x)$ as formal power series in $t$ with $C^\infty$ coefficients. The equations are obtained by writing
\begin{align}
&-h^2\Delta a^t+2(\nabla\psi_t+\ii tA)h\nabla a^t
+h[\Delta\psi_t+\ii t\operatorname{div}A]a^t
\notag\\
&\hspace{1.5cm}
+\bigl(V-(\nabla\psi_t+\ii tA)^2-E(h)\bigr)a^t
=\OO\bigl((|t|+|h|)^\infty\bigr).
\tag{2.41}
\end{align}
The first is the eikonal equation
\begin{equation}
V-(\nabla\psi_t+\ii tA)^2=\OO(|t|^\infty),
\tag{2.42}
\end{equation}
which was solved in (2.26)--(2.30).
The second equation, corresponding to the coefficient of $h$, consists in finding $E_1^t$ such that
\begin{equation}
2(\nabla\psi_t+\ii tA)\nabla a_0^t
+(\Delta\psi_t+\ii t\operatorname{div}A-E_1^t)a_0^t
=\OO(|t|^\infty),
\tag{2.43}
\end{equation}
with $a_0^t(0)=1$ for normalization.
A necessary condition is
\begin{equation}
E_1^t=(\Delta\psi_t+\ii t\operatorname{div}A)(0),
\tag{2.44}
\end{equation}
and the formula is indeed gauge invariant. This is a formal power series in $t$ because $\psi_t$ is. We used here $(\nabla\psi_t+\ii tA)(0)=0$.

\noindent\textbf{Exercise.}
Check the compatibility of (2.44), (2.28), and (2.20).

Expand (2.43) in powers of $t$ and look for $a_0^t$ in the form
\begin{equation}
a_{0,0}+\ii t a_{0,1}-t^2a_{0,2}+\cdots,
\qquad a_{0,j}\ \text{real},
\tag{2.45}
\end{equation}
and recall that
\begin{equation}
\psi^t=\psi_0+\ii t\psi_1-t^2\psi_2+\cdots,
\qquad \psi_j\ \text{real}.
\tag{2.46}
\end{equation}
One obtains
\begin{equation}
\begin{aligned}
2\nabla\psi_0\cdot\nabla a_{0,0}+(\Delta\psi_0-E_0^1)a_{0,0}&=0,\\
2\nabla\psi_0\cdot\nabla a_{0,1}
+(\Delta\psi_1+\operatorname{div}A-E_1^1)a_{0,0}&\\[-0.4ex]
\qquad{}+2(\nabla\psi_1+A)\cdot\nabla a_{0,0}&=0,\\
2\nabla\psi_0\cdot\nabla a_{0,j}
&=\varphi_j(x,a_{0,0},\ldots,a_{0,j-1},\\
&\hspace{2.0cm}\nabla a_{0,0},\ldots,\nabla a_{0,j-1})-E_1^j a_{0,0}.
\end{aligned}
\tag{2.47}
\end{equation}
The initial conditions are
\[
a_{0,0}(0)=1,\qquad a_{0,j}(0)=0.
\]
At each step $E_1^j$ is chosen so that the right-hand side vanishes at $0$ (this is summarized in (2.44)). These equations are solved term by term as in [HE--SJ]$_1$ (or more directly as in [HE]), giving a solution of (2.43).

\medskip
\noindent\textbf{Remark 2.11.}
The forms (2.45) and (2.46) are natural in view of Remark 2.9: if $u$ is a solution for $E^t$, then $\overline u$ is a solution for $E^{-t}$.

Finally consider the third equation (the subsequent ones are solved in the same way):
\begin{equation}
\begin{cases}
2(\nabla\psi_t+\ii tA)\nabla a_1^t
+(\Delta\psi_t+\ii t\operatorname{div}A-E_1^t)a_1^t
=\Delta a_0^t+E_2^t a_0^t,\\
a_1^t(0)=0.
\end{cases}
\tag{2.48}
\end{equation}
The solvability condition is
\begin{equation}
E_2^t=-(\Delta a_0^t)(0).
\tag{2.49}
\end{equation}
The equation is then solved in the same way as (2.43), via (2.47).

The preceding computations are summarized as follows.

\medskip
\noindent\textbf{Proposition 2.12 ($C^\infty$ case).}
Under the hypotheses of Proposition 2.4 and condition (2.40), there exists a symbol $\widetilde a^t(h)$ (a realization of a formal symbol in $t$ and $h$ with $C^\infty$ coefficients) and a normalization of the first eigenfunction $u_t(h)$ of $P_{t,\Omega}(h)$ such that
\[
u^t(h)(x)-h^{-n/4}\widetilde a^t\e^{-\widetilde\psi_t(x)/h}
=\OO(h^\infty)\e^{-\Ree\widetilde\psi_t(x)/h}
=\OO(h^\infty)\e^{-d_V(x)/h}
\]
in an open set $U$ of the type described in (2.36). Moreover $a_{0,0}(x)$ does not vanish on $\overline U$.

\section{Computation of the interaction matrix}

We saw in Section 2 how to treat the one-well problem. We now return to the determination of the spectrum in the case of several wells.
Let $I(h)\to0$ be an interval and assume that
\[
V^{-1}(0)=\bigcup_{j=1}^N U_j,
\]
where the $U_j$ are distinct compact wells of zero diameter for the Agmon distance. We retain the notation and assumptions of Section 1, in particular (1.3)--(1.10), and, in order to avoid difficulties when applying the results of Section 2, assume
\begin{equation}
m_j^t=1.
\tag{3.1}
\end{equation}
Thus there is a single eigenvalue $\mu_j^t(h)$ in $I(h)$ for the Dirichlet realization in $M_j(\eta)$, and it is easy to see that $\mu_j^t(h)$ depends continuously on $t$ (cf. [KA]$_2$). We obtain the analogue of Theorem 2.9 of [HE--SJ]$_1$.

\medskip
\noindent\textbf{Theorem 3.1.}
Under the hypotheses stated at the beginning of this section, the spectrum of $P^t(h)$ in $I(h)$ is given, modulo $\Ot(\e^{-2S_0/h})$ (that is, modulo $\OO(\e^{-2(S_0-\varepsilon(\eta))/h})$ with $\varepsilon(\eta)\to0$ as $\eta\to0$), by the matrix
\begin{equation}
\diag(\mu_j^t)+\widehat W_{j,k}^t,
\tag{3.2}
\end{equation}
where
\[
\widehat W_{j,k}^t=\frac12\left(W_{j,k}^t+\overline{W_{k,j}^t}\right)
\]
and
\begin{equation}
W_{j,k}^t
=h\int \chi_j\left[
\varphi_k^t\,\overline{\nabla_{tA}^h\varphi_j^t}
-\nabla_{tA}^h\varphi_k^t\,\overline{\varphi_j^t}
\right]\cdot\nabla\chi_k\dd x.
\tag{3.3}
\end{equation}
Here $\varphi_j^t$ is a normalized eigenfunction of $P_{M_j(\eta)}^t$ associated with $\mu_j$, $\chi_j$ is a $C^\infty$ function equal to $1$ on $M_j(2\eta)$ and supported in $M_j(3\eta/2)$, and
\[
\nabla_{tA}^h=h\nabla+\ii tA.
\]

\noindent\textbf{Proof.}
The proof is analogous to the case $t=0$ and relies on Green's formula.

\medskip
\noindent\textbf{Remark 3.2.}
\[
\widehat W_{j,j}^t=\Ot(\e^{-2S_0/h}),
\]
which follows immediately from Lemma 1.2 of [HE--SJ]$_1$.

\medskip
\noindent\textbf{Remark 3.3.}
The question is now whether the interaction matrix $\widehat W_{j,k}^t$, which is of order $\Ot(\e^{-d(U_j,U_k)/h})$, can be determined more explicitly. Let us spell out in our setting Remark (2.13) of [HE--SJ]$_1$.
Suppose $j\ne k$ are such that
\begin{equation}
d_V(U_j,U_k)=S_0,
\tag{3.4}
\end{equation}
and
\begin{equation}
\mu_j^t-\mu_k^t=\OO(\e^{-a/h})
\qquad\text{for some }a>0\text{ uniformly in }t.
\tag{3.5}
\end{equation}
Then
\begin{equation}
W_{jk}^t=\overline{W_{kj}^t}+\Ot(\e^{-S_0/h-a/h}),
\tag{3.6}
\end{equation}
with the same proof as when $t=0$. This holds in particular whenever
\begin{equation}
\mu_j^t=\mu_k^t,
\tag{3.7}
\end{equation}
which almost always follows from symmetry considerations.

\medskip
\noindent\textbf{Lemma 3.4.}
Suppose there is an isometry $\phi$ of $\R^n$ mapping $M_j$ onto $M_k$ such that
\begin{equation}
V\circ\phi=V,
\qquad
\phi_*A=A+\nabla\Theta,
\tag{3.8}
\end{equation}
where $\Theta$ is a $C^\infty$ function. Then $P^t_{M_j(\eta)}$ and $P^t_{M_k(\eta)}$ are unitarily equivalent. In particular, (3.7) holds.

\medskip
\noindent\textbf{Remark 3.5.}
Another potentially interesting situation is obtained by replacing (3.8) with
\begin{equation}
V\circ\phi=V,
\qquad
\phi_*A=-A+\nabla\Theta.
\tag{3.8$'$}
\end{equation}
Then there is a unitary $U$ such that the corresponding Dirichlet operators are intertwined after complex conjugation.

\medskip
\noindent\textbf{Remark 3.6.}
If (3.7) holds, then
\begin{equation}
W_{jk}^t=\overline{W_{kj}^t}.
\tag{3.9}
\end{equation}

\medskip
\noindent\textbf{Remark 3.7.}
Let $\Omega$ be a regular open set containing $U_j$ and contained in the complement of $U_k$, and let
\[
\Gamma=\partial\Omega\cap\{x:d(U_j,x)+d(U_k,x)\le S_0+a\}.
\]
Under hypothesis (3.5), Stokes' formula gives an expression analogous to formula (2.25) of [HE--SJ]$_1$:
\begin{align}
W_{jk}^t
={}&h\int_\Gamma\Bigg\{
\overline{\varphi_j^t}
\left(h\frac{\partial}{\partial n}+\ii tA\cdot n\right)\varphi_k^t
-\varphi_k^t\,
\overline{\left(h\frac{\partial}{\partial n}+\ii tA\cdot n\right)\varphi_j^t}
\Bigg\}\dd\nu_\Gamma
\notag\\
&+\Ot(\e^{-(S_0+a)/h}),
\tag{3.10}
\end{align}
where $n$ is the normal to $\Gamma$ oriented from $U_j$ toward $U_k$ and $\partial/\partial n$ is the normal derivative.

From now on assume that we can work with $I(h)=[0,Ch]$, where $C$ is chosen so that each Dirichlet realization $P^t_{M_j(\eta)}(h)$ associated with a well $U_j$, assumed to be a nondegenerate point well, has only one eigenvalue.
Assume that $V$ and $A$ are analytic in a neighborhood of
\[
d(x,U_j)+d(x,U_k)\le a+S_0.
\]
If the potentials are only $C^\infty$, restrict $t$ to
$|t|\le C\sqrt{h\log(1/h)}$.
Using the bottom-of-the-well B.K.W. constructions of Section 2, after a possible renormalization we may assume
\begin{equation}
\varphi_j^t(x,h)=h^{-n/4}a_j^t(x,h)\e^{-\psi_j^t(x)/h},
\tag{3.11}
\end{equation}
with
\begin{equation}
a_j^t(x,h)=\sum_{k=0}^\infty a_{jk}^t(x)h^k
\tag{3.12}
\end{equation}
in the sense of analytic symbols in the analytic category. The phase $\psi_j^t(x)$ is that constructed in Section 2 and $a_{j0}^t$ is elliptic. Recall that $\psi_j^0(x)=d_V(x,U_j)$.
This approximation is valid in a neighborhood of $\Gamma$, after possibly decreasing $\Gamma$ (that is, taking $a>0$ smaller) and $t_0>0$.

Substituting (3.11) and (3.12) into (3.10) gives, for some $\varepsilon_0>0$ and $t_0>0$, uniformly for $t\in[-t_0,t_0]$ and $h\in(0,h_0]$,
\begin{align}
W_{jk}^t
={}&h^{1-n/2}\int_\Gamma
\overline{a_j^t}a_k^t
\left[
\frac{\partial\psi_k^t}{\partial n}
+2\ii tA\cdot n
-\overline{\frac{\partial\psi_j^t}{\partial n}}
\right]
\e^{-\psi_{jk}^t/h}\dd\nu_\Gamma
\notag\\
&+h^{2-n/2}\int_\Gamma
\left[-a_k^t\overline{\frac{\partial a_j^t}{\partial n}}
+\overline{a_j^t}\frac{\partial a_k^t}{\partial n}\right]
\e^{-\psi_{jk}^t/h}\dd\nu_\Gamma
+\OO(\e^{-\varepsilon_0/h-S_0/h}),
\tag{3.13}
\end{align}
where, by definition,
\begin{equation}
\psi_{jk}^t(x)=\overline{\psi_j^t(x)}+\psi_k^t(x).
\tag{3.14}
\end{equation}
Recall that $\psi_{jk}^0(x)=d(x,U_j)+d(x,U_k)$.

\medskip
\noindent\textbf{Remark 3.8.}
In the $C^\infty$ case, the remainder $\OO(\e^{-\varepsilon_0/h-S_0/h})$ must be replaced by $\OO(h^\infty)\e^{-S_0/h}$ and one imposes $|t|\le C\sqrt{h\log(1/h)}$.

To measure more precisely the size of (3.13), note from (2.29) that
\begin{equation}
\Ree\psi_{jk}^t(x)\ge\Ree\psi_{jk}^0(x)=\psi_{jk}^0(x)\ge S_0.
\tag{3.15}
\end{equation}
It is important to know whether $\Ree\psi_{jk}^t(x)\ge S_0$ is strict for every $x\in\Gamma$ and to compute $\Ree\psi_{jk}^t(x)-S_0$ as a power series in $t$.
Return to the expansion
\begin{equation}
\psi_j^t(x)=\sum_{k=0}^\infty\psi_{j;k}(x)(\ii t)^k
\tag{3.16}
\end{equation}
and recall that $\psi_{j;2}(x)$ satisfies (2.27)$_2$:
\begin{equation}
2(\nabla\psi_{j;0})(x)\cdot\nabla\psi_{j;2}(x)
=-(\nabla\psi_{j;1}(x)+A)^2.
\tag{3.17}
\end{equation}
Fix a point $x$ for which there is a unique minimal geodesic $\gamma_j$ joining $x$ to $U_j$; after shrinking $\Gamma$ this is always the case for $x\in\Gamma$. Temporarily omit the index $j$.

\medskip
\noindent\textbf{Lemma 3.9.}
$\psi_2(x)=0$ if and only if $\nabla\psi_0\mathbin{\lrcorner}B$ vanishes on the minimal geodesic joining $x$ to the well $U$, where
\[
B=\sum_{j,k}\left(\frac{\partial A_j}{\partial x_k}-\frac{\partial A_k}{\partial x_j}\right)\dd x_j\wedge\dd x_k,
\qquad
\psi_0=d(x,U).
\]

\noindent\textbf{Proof of Lemma 3.9.}
Parametrize $\gamma$ on $(-\infty,0]$ as an integral curve of $\nabla\psi_0$. Then
\begin{equation}
\psi_2(x)=-\frac12\int_{-\infty}^0
|\nabla\psi_1+A|^2(\gamma(s))\dd s.
\tag{3.18}
\end{equation}
Hence clearly
\begin{equation}
\psi_2(x)=0
\quad\Longleftrightarrow\quad
(\nabla\psi_1+A)=0\quad\text{on the minimal geodesic }\gamma.
\tag{3.19}
\end{equation}
Since $\nabla\psi_1+A$ vanishes at the well, this is equivalent to
\begin{equation}
\psi_2(x)=0
\quad\Longleftrightarrow\quad
\frac{d}{ds}(\nabla\psi_1+A)(\gamma(s))=0,
\qquad s\in(-\infty,0].
\tag{3.20}
\end{equation}
Expanding the condition on the right gives
\begin{equation}
\sum_j\frac{\partial\psi_0}{\partial x_j}
\frac{\partial^2\psi_1}{\partial x_j\partial x_k}
+\sum_j\frac{\partial\psi_0}{\partial x_j}\frac{\partial A_k}{\partial x_j}=0
\quad\text{on }\gamma.
\tag{3.21}
\end{equation}
Recall also from (2.27)$_1$ that
\[
2\nabla\psi_0\cdot\nabla\psi_1=-2\nabla\psi_0\cdot A.
\]
Differentiating gives
\begin{equation}
\sum_j\frac{\partial^2\psi_0}{\partial x_j\partial x_k}
\left(\frac{\partial\psi_1}{\partial x_j}+A_j\right)
+\sum_j\frac{\partial\psi_0}{\partial x_j}
\frac{\partial^2\psi_1}{\partial x_j\partial x_k}
+\sum_j\frac{\partial\psi_0}{\partial x_j}\frac{\partial A_j}{\partial x_k}=0.
\tag{3.22}
\end{equation}
Thus (3.21) is equivalent to
\begin{equation}
\left.\sum_j\frac{\partial\psi_0}{\partial x_j}B_{jk}\right|_{\operatorname{Im}\gamma}
+\left.\sum_j\frac{\partial^2\psi_0}{\partial x_j\partial x_k}
\left(\frac{\partial\psi_1}{\partial x_j}+A_j\right)\right|_{\operatorname{Im}\gamma}=0.
\tag{3.23}
\end{equation}
It is then clear that $(\partial\psi_1/\partial x_j+A_j)=0$ on $\gamma$ implies $(\nabla\psi_0\mathbin{\lrcorner}B)=0$ there. Conversely, one obtains the system
\[
\begin{cases}
\displaystyle\frac{d}{ds}v_j(s)
=-\sum_k\frac{\partial^2\psi_0}{\partial x_j\partial x_k}(\gamma(s))v_k(s),\\[1.2ex]
v_j(-\infty)=0,
\end{cases}
\qquad s\in(-\infty,0],
\]
where $v_j(s)=(\nabla\psi_1+A)_j(\gamma(s))$. Its unique solution is $v_j(s)=0$, since the matrix $\partial^2\psi_0/\partial x_j\partial x_k(0)$ is nondegenerate positive. This proves Lemma 3.9.

\medskip
\noindent\textbf{Proposition 3.10 (one geodesic)} (cf. Theorem 6.6 of [HE--SJ]$_1$).
Suppose that there is only one minimal geodesic $\gamma_{jk}$ between $U_j$ and $U_k$, and that $\nabla\psi_{j;0}\mathbin{\lrcorner}B$ (equivalently $\nabla\psi_{k;0}\mathbin{\lrcorner}B$) is not identically zero on $\gamma_{jk}$. Then there exist $t_0>0$, $\alpha>0$ and $C>0$ such that, for $t\in[-t_0,t_0]$,
\begin{equation}
C\e^{-S_0/h}\e^{-\alpha t^2/h}\ge |W_{jk}^t|.
\tag{3.25}
\end{equation}

\noindent\textbf{Proof.}
If $x_{jk}=\gamma_{jk}\cap\Gamma$, then
\begin{align*}
\Ree\bigl(\overline{\psi_j^t}+\psi_k^t\bigr)(x_{jk})
={}&S_0+\frac12t^2\int_{-\infty}^0
(\nabla\psi_j^1+A)^2(\gamma_{jk}(s))\dd s\\
&+\frac12t^2\int_0^{+\infty}
(\nabla\psi_k^1+A)^2(\gamma_{jk}(s))\dd s
+\OO(t^4),
\end{align*}
which gives the result.

\medskip
\noindent\textbf{Remark 3.11.}
One may of course ask for a lower bound. This is more delicate than in the case $t=0$. Indeed, in formula (3.13), in a neighborhood of $x_{jk}$,
\[
\overline{a_j^t}a_k^t\left(\frac{\partial\psi_k^t}{\partial n}
+2\ii t A\cdot n-
\overline{\frac{\partial\psi_j^t}{\partial n}}\right)
\]
remains, for small $t$, close to
$-2a_j^0(x_{jk})a_k^0(x_{jk})\sqrt{V(x_{jk})}$, but an imaginary part appears in the phase and can cause cancellations in the integral, making the splitting smaller. In particular, the argument of $W_{jk}^t$, which is $\pi$ for $t=0$, may take other values. Moreover, if there are several geodesics, compensation phenomena can occur and may even lead to cancellation of the tunneling effect. We shall now try to analyze these phenomena by looking for cases in which one can obtain an asymptotic equivalent for $W_{jk}^t$.

\medskip
\noindent\textbf{Proposition 3.12 (one geodesic).}
Suppose there is only one minimal geodesic $\gamma_{jk}$, that $\Gamma$ intersects $\gamma_{jk}$ transversely at $x_{jk}$, and that the restriction to $\Gamma$ of
\[
 d(x,U_j)+d(x,U_k)
\]
vanishes exactly to order two at $x_{jk}$. Then there exist $a>0$ and $t_0>0$ such that, for $t\in[-t_0,t_0]$,
\begin{equation}
W_{jk}^t
=h\,b_{jk}^t(h)\e^{-S_{jk}^t/h}
+\OO\!\left(\e^{-(S_0+a)/h}\right),
\tag{3.26}
\end{equation}
where $b_{jk}^t(h)$ is an elliptic analytic symbol, and the expansion of $S_{jk}^t$ through order two can be computed; see (3.32) and (3.33). In particular, there is $C>0$ such that
\[
 |W_{jk}^t|\ge Ch\e^{-\Ree S_{jk}^t/h}.
\]

\begin{quote}\small
\textit{Translator's note.} In the scan, the displayed formula numbered (3.26) is not visibly printed in the statement of Proposition 3.12. The formula above is reconstructed from the immediately following discussion and from Remark 3.17, which gives its $C^\infty$ counterpart explicitly.
\end{quote}

\noindent\textbf{Proof of Proposition 3.18 [as printed; Proposition 3.12 is meant].}
This is an application of the analytic stationary-phase theorem. The phase is the function $\psi_{jk}^t(x)$, regarded as an analytic function on $\Gamma$ and extended to a complex neighborhood, and the amplitude is likewise extended:
\begin{align*}
x\longmapsto c_{jk}^t(h,x)
={}&\overline{a_j^t(h,x)}a_k^t(h,x)
\left[
\frac{\partial\psi_k^t}{\partial n}
+2\ii tA\cdot n
-\overline{\frac{\partial\psi_j^t}{\partial n}}
\right]\\
&+h\left[-a_k^t\overline{\frac{\partial a_j^t}{\partial n}}
+\overline{a_j^t}\frac{\partial a_k^t}{\partial n}\right].
\end{align*}
With this notation,
\[
W_{jk}^t=\int_\Gamma \e^{-\psi_{jk}^t(x)/h}
 c_{jk}^t(h,x)\dd\nu_\Gamma,
\]
and $\Ree\psi_{jk}^t\ge S_0+a$ on $\partial\Gamma$ for $t\in[-t_0,t_0]$, with $t_0$ small enough. We deform the contour so that it passes through the critical point $x_{jk}^t$ of $\psi_{jk}^t$ on $\Gamma^{\C}$, with $x_{jk}^0=x_{jk}$. Its existence follows from the implicit-function theorem, thanks to the nondegeneracy assumption on the Hessian of the restriction to $\Gamma$ of
$d(x,U_j)+d(x,U_k)$.

An application of Theorem (2.8) of [SJ] gives (3.26) with
\begin{equation}
S_{jk}^t=\psi_{jk}^t(x_{jk}^t).
\tag{3.27}
\end{equation}
First observe that
\begin{equation}
\Argg b_{jk}^t(h)=\pi+\OO(|t|+|h|),
\tag{3.28}
\end{equation}
and
\begin{equation}
\psi_{jk}^t(x_{jk}^t)=S_0+\OO(|t|).
\tag{3.29}
\end{equation}
We now determine $S_{jk}^t$ more precisely by computing the critical point and then working modulo $\OO(t^3)$.

\paragraph{Determination of the critical points.}
As a preliminary calculation, consider a family of holomorphic functions $\rho_t$ in a complex neighborhood of $0$, depending holomorphically on a complex parameter $t$, such that
\[
\rho_t(x)=\rho_0(x)+\ii t\rho_1(x)+t^2\rho_2(x)+\OO(t^3),
\]
where $\rho_1$ is real for real $x$, $\rho_0(0)=0$, $\rho_0'(0)=0$, $\rho_0''(0)$ is positive definite, and $\Ree\rho_0\ge0$. Near $0$, the function $x\mapsto\rho_t(x)$ has a critical point
\begin{equation}
x_t=-\ii t\,\rho_0''(0)^{-1}\rho_1'(0)+\OO(t^2).
\tag{3.30}
\end{equation}
At that critical point, for real $t$,
\begin{align}
\Ree\rho_t(x_t)
&=\left[-\rho_2(0)
+\frac12\bigl(\rho_0''(0)^{-1}\rho_1'(0)\mid\rho_1'(0)\bigr)\right]t^2
+\OO(t^4),\notag\\
\Imm\rho_t(x_t)&=t\rho_1(0)+\OO(t^3).
\tag{3.31}
\end{align}
Applying this elementary calculation to $\psi_{jk}^t$ restricted to $\Gamma$ gives
\begin{align}
\Ree S_{jk}^t
={}&S_0+t^2\Bigg[-\psi_k^2(x_{jk})-\psi_j^2(x_{jk})\notag\\
&\qquad +\frac12\left\langle
\nabla\psi_k^1-\nabla\psi_j^1,
\bigl(\rho_0''(x_{jk})\bigr)^{-1}
(\nabla\psi_k^1-\nabla\psi_j^1)
\right\rangle\Bigg]
+\OO(t^4),
\tag{3.32}
\end{align}
and
\begin{equation}
\Imm S_{jk}^t
=t\bigl(\psi_k^1(x_{jk})-\psi_j^1(x_{jk})\bigr)+\OO(t^3).
\tag{3.33}
\end{equation}
Here $\rho_0''(x_{jk})$ denotes the Hessian of the restriction to $\Gamma$ of the unperturbed phase. Formula (3.32) is meaningful because, by (2.27)$_1$, the vectors $\nabla\psi_k^1+A$ and $\nabla\psi_j^1+A$ are orthogonal to $\gamma_{jk}$ at $x_{jk}$; hence their difference is tangent to $\Gamma$.
This proves the proposition.

\medskip
\noindent\textbf{Proposition 3.13.}
Under the hypotheses of Proposition 3.12, there is $\alpha>0$ such that
\[
\Ree S_{jk}^t\ge S_0+\alpha t^2
\qquad (|t|\le t_0\text{ sufficiently small})
\]
if and only if $\nabla\psi_0\mathbin{\lrcorner}B$ is not identically zero on the minimal geodesic joining $U_j$ to $U_k$.

Indeed, the coefficient of $t^2$ in $S_{jk}^t$ is a sum of three positive terms, and vanishing of the first two forces the third to vanish as well.

A second natural question is to interpret the first coefficients in powers of $t$ more intrinsically and more precisely. In particular, one can expect to see directly that they are independent of the choice of the point $x_{jk}$. The independence of $S_{jk}^t$ from $x_{jk}$ is clear, for example, from (3.26) and
\[
S_{jk}^t=\lim_{h\to0}(-h)\log(-W_{jk}^t),
\]
where the principal determination of the logarithm is used.

\medskip
\noindent\textbf{Remark 3.14 (calculation of $\Imm S_{jk}^t$).}
We compute the coefficient of $t$; recall that $x_{jk}$ is an arbitrary point on the minimal geodesic $\gamma_{jk}$ between $U_j$ and $U_k$. By (2.27)$_1$, and with the normalization (2.26)$_3$,
\begin{align*}
\psi_j^1(\gamma_{jk}(r))
&=-\int_{-\infty}^{r}
\dot\gamma_{jk}(s)\cdot A(\gamma_{jk}(s))\dd s,\\
\psi_k^1(\gamma_{jk}(r))
&=-\int_{-\infty}^{-r}
\dot\gamma_{jk}(-s)\cdot A(\gamma_{jk}(-s))\dd s\\
&=\int_{r}^{+\infty}
\dot\gamma_{jk}(s)\cdot A(\gamma_{jk}(s))\dd s.
\end{align*}
We obtain:

\medskip
\noindent\textbf{Lemma 3.15.}
\begin{equation}
\Imm S_{jk}^t
=t\,\Circ(A,\gamma_{jk})+\OO(t^3),
\tag{3.34}
\end{equation}
where $\Circ(A,\gamma_{jk})$ is the circulation of $A$ along $\gamma_{jk}$.

This formula is not gauge invariant, since the normalization (2.26)$_3$ is itself not invariant under a gauge transformation. It would also be interesting to have an elegant expression for
\[
\lim_{t\to0}\frac{\Ree S_{jk}^t-S_0}{t^2}.
\]

\medskip
\noindent\textbf{Remark 3.16 (relation with an action integral).}
Let us formally reconsider the preceding calculations in order to recover the action integrals used by physicists (cf. [WI]). Let
\[
\psi_j^t\text{ be the phase relative to }U_j,
\qquad \phi_j^t=\ii\psi_j^t,
\]
and
\[
\psi_k^t\text{ be the phase relative to }U_k,
\qquad \phi_k^t=\ii\psi_k^t,
\qquad \widetilde\phi_k^t(x)=\overline{\phi_k^t(\bar x)}.
\]
Set
\[
\Lambda_{\phi_j^t}=\{(x,\xi):\xi=\phi_j^{t\prime}(x)\},
\qquad
\Lambda_{\widetilde\phi_k^t}=\{(x,\xi):\xi=\widetilde\phi_k^{t\prime}(x)\}.
\]
If
\[
p^t(x,\xi)=(\xi-tA(x))^2+V(x),
\]
then
\begin{equation}
p^t(x,\phi_j^{t\prime}(x))=0,
\qquad
p^t(x,\widetilde\phi_k^{t\prime}(x))=0.
\tag{3.35}
\end{equation}
Finding the critical point $x_{jk}^t\in\Gamma^{\C}$ of
$\phi_j^t-\widetilde\phi_k^t$ restricted to $\Gamma^{\C}$ amounts to imposing
\begin{equation}
\nabla_{\Gamma^{\C}}\phi_j^t
=\nabla_{\Gamma^{\C}}\widetilde\phi_k^t.
\tag{3.36}
\end{equation}
Thus the critical point lies at the intersection of
$\Lambda_{\phi_j^t}$ and $\Lambda_{\widetilde\phi_k^t}$, which intersect along a complex curve. Choose a curve $\widehat\gamma_t$ in this intersection joining $U_j$ to $U_k$; for instance one may use the curve traced out by the points $x_{jk}^t$ when $\Gamma$ varies. Consider the action integral
\begin{equation}
I(t)=\int_{\widehat\gamma_t}\xi\,\dd x.
\tag{3.37}
\end{equation}
Clearly,
\begin{align*}
I(t)
&=\int_{-\infty}^{0}\frac{d}{ds}
\phi_j^t(\widehat\gamma_t(s))\dd s
 +\int_{0}^{+\infty}\frac{d}{ds}
\widetilde\phi_k^t(\widehat\gamma_t(s))\dd s\\
&=\phi_j^t(x_{jk}^t)-\widetilde\phi_k^t(x_{jk}^t)
=\ii\bigl(\psi_j^t(x_{jk}^t)+\psi_k^t(x_{jk}^t)\bigr),
\end{align*}
whence
\begin{equation}
I(t)=\ii S_{jk}^t.
\tag{3.38}
\end{equation}
By Stokes' formula,
\begin{equation}
\dot I(t)=\int_{-\infty}^{+\infty}
\sigma\!\left(\frac{\partial\widehat\gamma_t}{\partial t},H_{p^t}\right)\dd s,
\tag{3.39}
\end{equation}
and in particular
\begin{align}
\dot I(0)
&=\psi_k^1(x_{jk})-\psi_j^1(x_{jk})\notag\\
&=\int_{-\infty}^{+\infty}
\sigma\!\left(
\left.\frac{\partial\widehat\gamma_t}{\partial t}\right|_{t=0},H_{p^0}
\right)\dd s\notag\\
&=\int_{-\infty}^{+\infty}
\left\langle \dd p^0,
\left.\frac{\partial\widehat\gamma_t}{\partial t}\right|_{t=0}
\right\rangle(\widehat\gamma_0(s))\dd s.
\tag{3.40}
\end{align}
Since $p^t(\widehat\gamma_t(s))=0$, we have
\[
\frac{\partial p^t}{\partial t}(\widehat\gamma_t)
+\dd p^t\cdot\frac{\partial\widehat\gamma_t}{\partial t}=0.
\]
Consequently,
\begin{align}
\dot I(0)
&=-\int_{-\infty}^{+\infty}
\left.\frac{\partial p^t}{\partial t}\right|_{t=0}
(\widehat\gamma_0(s))\dd s,\notag\\
\dot I(0)
&=\int_{-\infty}^{+\infty}
\dot\gamma_{jk}(s)\cdot A(\gamma_{jk}(s))\dd s.
\tag{3.41}
\end{align}
Thus (3.34) is recovered, at least formally.

\medskip
\noindent\textbf{Remark 3.17.}
When $V$ and $A$ are only $C^\infty$, formula (3.26) can still be given a meaning provided one imposes
\[
|t|\le C\sqrt{h}\sqrt{\log(1/h)}.
\]
Then
\[
W_{jk}^t=h\,b_{jk}^t(h)\e^{-S_{jk}^t/h}
+\OO(h^\infty)\e^{-S_0/h},
\]
where $b_{jk}^t(h)$ and $S_{jk}^t$ are realizations of formal series in $(t,h)$. It suffices to apply a formal stationary-phase theorem.

\section{Applications to the Study of the ``Splitting'' of Eigenvalues}

We return to the hypotheses considered in (3.5). We assume that there are two nondegenerate wells $U_1$ and $U_2$, and that there is an isometry $\phi$ of $\R^n$ such that
\begin{align}
\phi^2&=I, \tag{4.1}\\
V\circ\phi&=V, \tag{4.2}\\
\phi_*A&=A+\nabla\Theta
\qquad\text{for some }\Theta\in C^\infty(\R^n), \tag{4.3}\\
\phi(U_1)&=U_2. \tag{4.4}
\end{align}
Observe first that by replacing $A$ with
$\frac12(\phi_*A+A)$, which is gauge equivalent since
\[
\frac12(\phi_*A+A)-A=\nabla(\Theta/2),
\]
we may replace (4.3) by
\[
\phi_*A=A. \tag{4.3$'$}
\]

We now follow closely the calculations of [HE--SJ]$_2$. In the ``adapted'' basis, the $2\times2$ matrix that determines the first two eigenvalues near the first eigenvalue of the harmonic oscillator (the bottom of the well) depends continuously on $t$, is Hermitian, and commutes with the matrix
$\left(\begin{smallmatrix}0&1\\1&0\end{smallmatrix}\right)$ of the natural action of $\phi$ in the eigenspace under consideration. Thus it has the form
\begin{equation}
\begin{pmatrix}
\lambda(t,h)&\alpha(t,h)\\
\alpha(t,h)&\lambda(t,h)
\end{pmatrix},
\tag{4.5}
\end{equation}
with, in view of the preceding remarks and the calculations of Section~3,
\begin{align}
\alpha(t,h)&\in\R\quad\text{and depends continuously on }t,
\tag{4.6}\\
\lambda(t,h)&=\lambda_D^t(h)+\widetilde\OO(\e^{-2S_0/h})
\qquad (t\in[-t_0,t_0]),
\tag{4.7}\\
\alpha(t,h)&=W_{12}^t+\OO(\e^{-(S_0+a)/h}),\qquad a>0.
\tag{4.8}
\end{align}
Here $W_{12}^t$ is the interaction coefficient introduced in Section~3 (cf. (3.3), (3.9), and (3.10)), defined from the functions
\begin{equation}
\begin{aligned}
\varphi_1^t &: \text{ normalized eigenfunction of the Dirichlet problem in }M_1,\\
&\text{ associated with }\lambda_D^t(h),\\
\varphi_2^t&=\varphi_1^t\circ\phi^{-1}.
\end{aligned}
\tag{4.9}
\end{equation}

In general, three cases may occur:
\[
\text{(a) }\alpha(t,h)<0,
\qquad
\text{(b) }\alpha(t,h)=0,
\qquad
\text{(c) }\alpha(t,h)>0.
\]
Case (a) is the one occurring for $t=0$ (cf. [HE--SJ]$_1$): the smaller eigenvalue is $\lambda(t,h)+\alpha(t,h)$ and corresponds, in the adapted basis, to a symmetric state with coordinates $(1,1)$. Case (b) would correspond to a double eigenvalue, which is excluded a priori without a magnetic field, but is a priori possible in the magnetic case (cf. [A--H--S]). Case (c) is also specific to the magnetic field (if it occurs): the lower eigenvalue would correspond to an antisymmetric eigenfunction for $\phi$, vanishing on the fixed-point set of the isometry. The occurrence of case (c) implies the occurrence of case (b), by continuity in $t$ and the fact that $\alpha(0,h)<0$. We shall now show by examples that all of these situations can indeed occur under suitable assumptions.

\subsection*{The case of one geodesic}
Assume there is a unique nondegenerate minimal geodesic $\gamma_{12}$ between the two wells. Necessarily $\phi(\gamma_{12})=\gamma_{12}$. Since $\phi$ interchanges the two wells, it necessarily has a fixed point on $\gamma_{12}$. As $\phi$ is an isometry of $\R^n$, it is represented by an orthogonal matrix (whose square is $1$, since $\phi^2=1$). After changing coordinates, we may suppose that $0\in\gamma_{12}$ and take for $\Gamma$ an open subset of a $\phi$-invariant hyperplane $\Sigma$ transverse to $\gamma_{12}$.

This case brings nothing essentially new. Proposition 3.12 gives an exponential lower bound of the form
\[
|W_{12}^t|\gtrsim h\e^{-\Ree S_{12}^t/h},
\qquad W_{12}^0<0.
\]
Moreover $\alpha(t,h)<0$ for $t\in[-t_0,t_0]$ if $t_0$ is small enough. Thus the only observed effect is a decrease of the tunneling effect caused by the magnetic field when $\Ree S_{12}^t>S_0$.

\subsection*{The case of two geodesics}
Assume there are two nondegenerate minimal geodesics $\gamma_{12}$ and $\widetilde\gamma_{12}$ between the two wells, and that
\[
\phi(\gamma_{12})=\widetilde\gamma_{12}.
\]
We restrict ourselves to the case $n=2$ and take $\phi$ to be the symmetry about the origin,
\[
\phi=\begin{pmatrix}-1&0\\0&-1\end{pmatrix}.
\]
Suppose
\[
U_1=(-a_1,0),\qquad U_2=(a_1,0),
\]
and take $\Sigma=\{x_1=0\}$.

\begin{center}
\begin{tikzpicture}[scale=0.9]
  \draw[->] (-4.2,0)--(4.2,0) node[right] {$x_1$};
  \draw[->] (0,-2.2)--(0,2.2) node[above] {$x_2$};
  \fill (-3,0) circle (1.5pt) node[below] {$U_1$};
  \fill (3,0) circle (1.5pt) node[below] {$U_2$};
  \draw (-3,0) .. controls (-1.8,1.25) and (1.8,1.25) .. (3,0);
  \draw (-3,0) .. controls (-1.8,-1.25) and (1.8,-1.25) .. (3,0);
  \node at (1.4,1.15) {$\gamma_{12}$};
  \node at (1.5,-1.15) {$\widetilde\gamma_{12}$};
  \draw[dashed] (-0.25,1.0)--(0.25,1.0);
  \node[left] at (-0.05,1.0) {$a_2$};
  \draw[dashed] (-0.25,-1.0)--(0.25,-1.0);
  \node[left] at (-0.05,-1.0) {$-a_2$};
\end{tikzpicture}
\end{center}

Modulo $\OO(\e^{-(S_0+a)/h})$, $W_{12}^t$ is the sum of two contributions, conjugate to each other:
\begin{equation}
W_{12}^t=W_{12}^{t,(\gamma)}+W_{12}^{t,(\widetilde\gamma)},
\tag{4.10}
\end{equation}
with
\begin{equation}
\overline{W_{12}^{t,(\gamma)}}=W_{12}^{t,(\widetilde\gamma)},
\tag{4.11}
\end{equation}
as a consequence of the reality of $W_{12}^t$. Moreover the contribution associated with $\gamma_{12}$ is given by
\begin{align}
W_{12}^{t,(\gamma)}
={}&h^2\int_{|x_2-a_2|<\varepsilon}
\Bigg[
\overline{\varphi_1^t(0,x_2)}
\frac{\partial\varphi_1^t}{\partial x_1}(0,-x_2)
+\varphi_1^t(0,-x_2)
\overline{\frac{\partial\varphi_1^t}{\partial x_1}(0,x_2)}
\Bigg]\dd x_2\notag\\
&+2\ii h t\int_{|x_2-a_2|<\varepsilon}
A_1(0,x_2)\,
\overline{\varphi_1^t(0,x_2)}\varphi_1^t(0,-x_2)\dd x_2.
\tag{4.12}
\end{align}
Thus $W_{12}^t$ is real and may be written
\begin{equation}
W_{12}^t
=2\bigl|W_{12}^{t,(\gamma)}\bigr|
\cos\bigl(\Argg W_{12}^{t,(\gamma)}\bigr).
\tag{4.13}
\end{equation}
The contribution $W_{12}^{t,(\gamma)}$ can be calculated using Proposition 3.12, but what matters most here is its sign and hence its argument. Indeed, $\alpha(t,h)$ has the same sign as $W_{12}^t$ as soon as, for some $\varepsilon_0>0$ independent of $h$,
\[
\frac{\pi}{2}+\varepsilon_0
<\Argg W_{12}^{t,(\gamma)}
<\frac{3\pi}{2}-\varepsilon_0.
\]
Because of the lower bound on $|W_{12}^{t,(\gamma)}|$, the remainder cannot change the sign.

By Proposition 3.12,
\begin{equation}
\Argg W_{12}^{t,(\gamma)}
=\frac{\Imm S_{12}^{t,(\gamma)}}{h}
+\OO(|h|+|t|),
\tag{4.14}
\end{equation}
and, by (3.34),
\begin{equation}
\Imm S_{12}^{t,(\gamma)}
=t\,\Circ(A,\gamma_{12})+\OO(t^3).
\tag{4.15}
\end{equation}
The symmetries imply
\begin{equation}
\Circ(A,\gamma_{12})=-\Circ(A,\widetilde\gamma_{12}).
\tag{4.16}
\end{equation}
Denoting by $\Sigma(\gamma,\widetilde\gamma)$ the region enclosed by $\gamma$ and $\widetilde\gamma$, we obtain
\begin{equation}
\Circ(A,\gamma)
=\frac12\int_{\Sigma(\gamma,\widetilde\gamma)}B,
\qquad
B=\left(
\frac{\partial A_1}{\partial x_2}-
\frac{\partial A_2}{\partial x_1}
\right)\dd x_1\wedge\dd x_2.
\tag{4.17}
\end{equation}
Thus, as soon as
\begin{equation}
\int_{\Sigma(\gamma,\widetilde\gamma)}B\ne0,
\tag{4.18}
\end{equation}
we can, by varying $(t,h)$ with $0<|t|<t_0$ ($t_0$ depending on $h$ in the $C^\infty$ case) and $h\in(0,h_0]$, obtain any desired value of the argument of $W_{12}^{t,(\gamma)}$. In particular, for $t$ near
\begin{equation}
t_k(h)=\frac{2(k+\tfrac12)\pi h}
{\displaystyle\int_{\Sigma(\gamma,\widetilde\gamma)}B},
\qquad k\in\Z,
\tag{4.19}
\end{equation}
there is a double eigenvalue.

\medskip
\noindent\textbf{Remark 4.1.}
This effect seems to be new, but it is of the same type as that appearing in the non-rigorous article of Aharonov and Bohm [AH--BO].

\section*{References}
\begin{description}[leftmargin=3.2cm,style=nextline]
\item[{[AH--BO]}] Y. Aharonov and D. Bohm, \emph{Significance of Electromagnetic Potentials in the Quantum Theory}, The Physical Review, vol.~115 (3), Aug. 1959.

\item[{[A--H--S]}] J. Avron, I. Herbst, and B. Simon, \emph{Schrödinger Operators with Magnetic Fields}.\\
{}[1] \emph{General Interactions}, Duke Math. J., vol.~45 (4), Dec. 1978.\\
{}[2] \emph{Separation of Center of Mass in Homogeneous Magnetic Fields}, Ann. Physics 114, pp.~431--451 (1978).\\
{}[3] \emph{Atoms in Homogeneous Magnetic Fields}, Comm. Math. Phys. 79, pp.~529--572 (1981).

\item[{[C--S--S]}] J.-M. Combes, R. Schrader, and R. Seiler, \emph{Classical Bounds and Limits for Energy Distributions of Hamilton Operators in Electromagnetic Fields}, Ann. Physics 111, pp.~1--18 (1978).

\item[{[HE]}] B. Helffer, \emph{An introduction to Semi Classical Analysis for the Schrödinger Operator}. Cours en Chine. (A paraître Lecture Notes in Mathematics 1988).

\item[{[HE--RO]}] B. Helffer and D. Robert, \emph{Calcul Fonctionnel Admissible et Applications}, J. Funct. Anal. 53 (3) (1983), pp.~246--268.

\item[{[HE--SJ]}] B. Helffer and J. Sjöstrand, \emph{Multiple Wells in the Semi-Classical Limit}.\\
{}[1] Comm. in P.D.E. 9 (4) (1984), pp.~337--408.\\
{}[2] \emph{Interaction Moléculaire, Symétries, perturbation}, Ann. Inst. H. Poincaré, vol.~42 (2) (1985), pp.~127--212.

\item[{[HE--MA--RO]}] B. Helffer, A. Martinez, and D. Robert, \emph{Ergodicité et Limite Semi-classique} (Comm. Math. Phys.).

\item[{[HO]}] L. Hörmander, \emph{The Analysis of Linear Partial Differential Equations}, Tome 3, Chap. XXII, Springer.

\item[{[HU]}] W. Hunziker, \emph{Schrödinger Operators with Electric or Magnetic Fields}, Proc. Int. Conf. in Math. Phys., Lausanne (1980).

\item[{[IV]}] I. V. Ivrii, Note au C.R.A.S. 1986, 302 Sér. I (13), pp.~467--471.

\item[{[KA]}] T. Kato.\\
{}[1] Israel J. Math. 13 (1972), pp.~125--174.\\
{}[2] \emph{Perturbation Theory for Linear Operators}, Grundlehren der Mathematischen Wissenschaften 132, Springer Verlag.

\item[{[ME--SJ]}] A. Menikoff and J. Sjöstrand, \emph{On the Eigenvalues of a class of hypoelliptic Operators}, Math. Ann. 235, pp.~55--85 (1978).

\item[{[MO]}] A. Mohamed.

\item[{[SI]}] B. Simon.\\
{}[1] Indiana Univ. Math. J. 26, pp.~1067--1073 (1977).\\
{}[2] \emph{Semi-Classical Analysis of Low Lying Eigenvalues I}, Ann. Inst. H. Poincaré, t.~38 (1983), pp.~295--307.\\
{}[3] \emph{Semi-Classical Analysis of Low Lying Eigenvalues II. Tunneling}, Ann. of Math. (1984), vol.~20, pp.~89--118.

\item[{[SJ]}] J. Sjöstrand.\\
{}[1] \emph{Singularités analytiques microlocales}, Astérisque 95 (1982).\\
{}[2] \emph{Analytic Wave front Sets and Operators with Multiple Characteristics}, Hokkaido Math. Journal, vol.~XII (3), pp.~392--433.

\item[{[WI]}] M. Wilkinson, \emph{Narrowly avoided crossings}, preprint, April 1986.
\end{description}

\bigskip
\noindent
École Polytechnique\\
Centre de Mathématiques\\
Plateau de Palaiseau\\
91128 Palaiseau Cedex\\
France

\medskip
\noindent
Département de Mathématiques\\
Université Paris-Sud\\
91405 Orsay Cedex\\
France

\end{document}
